% Kandydat restrukturyzacji; nie jest aktywnym rozdziałem Atlasu.

\section{Splot na grupach lokalnie zwartych}
\label{sec:dt-haar-convolution}
\providecommand{\atlasref}[2]{#2}

Konstrukcję i jednoznaczność miary Haara pozostawiamy
w teorii miary. W tym rozdziale korzystamy z wyników
ME-23 dla lokalnie zwartej grupy Hausdorffa $G$
z lewostronną miarą Haara $\mu$. Nie zakładamy
$\sigma$-skończoności $G$.
Przyjmujemy konwencję funkcji modularnej
$\mu(Ea)=\Delta(a)\mu(E)$, a więc
\[
 \int_G h(x^{-1})\,\mathrm d\mu(x)
 =\int_G h(x)\Delta(x)^{-1}\,\mathrm d\mu(x)
\]
dla funkcji nieujemnych lub w przypadku bezwzględnej
zbieżności odpowiednich całek. Są to wyniki
\atlasref{ME-23::thm:haar-modular-23}{ME-23: modularność}
oraz \atlasref{ME-23::cor:haar-integral-inversion-23}{ME-23: inwersja}.
Por. także \atlasref{ME-23::it:def:haar-convolution}{ME-23: splot na grupie}.

\subsection{Splot funkcji ciągłych o zwartym nośniku}
\begin{definicja}[Splot na grupie]
\label{def:dt-haar-convolution}
Niech $G$ będzie lokalnie zwartą grupą Hausdorffa,
$\mu$ jej lewostronną miarą Haara, a
$\mathbb K=\mathbb R$ lub $\mathbb C$.
Dla $f,g\in C_c(G;\mathbb K)$ definiujemy
\[
 (f*g)(x)=\int_G f(y)g(y^{-1}x)\,\mathrm d\mu(y).
\]
Dla ustalonego $x$ funkcja podcałkowa jest ciągła
i znika poza zwartym nośnikiem $f$, więc całka
istnieje i jest skończona. Na grupie addytywnej
$\mathbb R^d$, po wyborze normalizacji $\mu=\lambda_d$,
otrzymujemy zwykły splot
$(f*g)(x)=\int_{\mathbb R^d}f(y)g(x-y)
\,\mathrm d\lambda_d(y)$.
Dla innej miary Haara $\mu=c\lambda_d$, $c>0$,
splot jest $c$ razy większy od splotu względem
$\lambda_d$. Wynika to z jednoznaczności miary Haara
(twierdzenie \atlasref{ME-23::thm:haar-unique-23}{ME-23: jednoznaczność miary Haara}) i liniowości
całki względem miary.
\end{definicja}

\begin{twierdzenie}[Podstawowe własności splotu]
\label{thm:dt-haar-convolution}
Dla $f,g,h\in C_c(G;\mathbb K)$ splot jest
dwuliniowy, $f*g\in C_c(G;\mathbb K)$ oraz
$\operatorname{supp}(f*g)\subseteq
(\operatorname{supp}f)(\operatorname{supp}g)$.
Ponadto
\[
 \int_G(f*g)\,\mathrm d\mu
 =\left(\int_G f\,\mathrm d\mu\right)
  \left(\int_G g\,\mathrm d\mu\right),
 \qquad
 \|f*g\|_1\le\|f\|_1\|g\|_1,
\]
\[
 \|f*g\|_\infty\le\|f\|_1\|g\|_\infty,\qquad
 (f*g)*h=f*(g*h).
\]
\end{twierdzenie}
\begin{proof}
Dwuliniowość wynika z liniowości całki.
Oznaczmy $K=\operatorname{supp}f$,
$L=\operatorname{supp}g$ i $M=\operatorname{supp}h$.
Jeżeli $f(y)g(y^{-1}x)\ne0$, to $y\in K$,
$y^{-1}x\in L$, a więc $x\in KL$. Zbiór $KL$
jest zwarty jako obraz $K\times L$ przez mnożenie,
a zatem domknięty. Stąd nośnik splotu zawiera się
w $KL$. Dla $a\to e$ mamy
\[
 |(f*g)(xa)-(f*g)(x)|
 \le\|f\|_1\sup_{z\in G}|g(za)-g(z)|\longrightarrow0
\]
na mocy lematu \atlasref{ME-23::lem:haar-uniform-23}{ME-23: jednostajna ciągłość translacji},
który w przypadku zespolonym stosujemy do części
rzeczywistej i urojonej. Ponieważ każde $x'$
ma postać $xa$ z $a=x^{-1}x'$, otrzymujemy
ciągłość splotu w każdym $x$.
Nierówność z normą supremum wynika bezpośrednio
z oszacowania $|g(y^{-1}x)|\le\|g\|_\infty$.

Jądro $(y,x)\mapsto f(y)g(y^{-1}x)$ jest ciągłe,
o nośniku zawartym w zwartym $K\times KL$.
Z lematu \atlasref{ME-23::lem:haar-fubini-compact-23}{ME-23: Fubini dla zwartych nośników}
i lewostronnej niezmienniczości mamy
\[
 \int_G(f*g)(x)\,\mathrm d\mu(x)
 =\int_G f(y)\left(\int_G g(y^{-1}x)
                 \,\mathrm d\mu(x)\right)\mathrm d\mu(y)
 =\int_G f\,\mathrm d\mu\int_G g\,\mathrm d\mu.
\]
Punktowa nierówność $|f*g|\le |f|*|g|$ oraz
powyższy wzór dla funkcji nieujemnych dają
nierówność norm $L^1$.

Dla ustalonego $x$ jądro
$(y,z)\mapsto f(y)g(y^{-1}z)h(z^{-1}x)$
jest ciągłe i znika poza $K\times KL$.
Możemy zatem zamienić kolejność całkowania
tym samym lematem Fubiniego. Po podstawieniu
$z=yu$ w całce względem $z$ dostajemy
\begin{align*}
 ((f*g)*h)(x)
 &=\int_G\int_G f(y)g(y^{-1}z)h(z^{-1}x)
       \,\mathrm d\mu(y)\,\mathrm d\mu(z)\\
 &=\int_G f(y)\int_G g(u)h(u^{-1}y^{-1}x)
       \,\mathrm d\mu(u)\,\mathrm d\mu(y)\\
 &= (f*(g*h))(x).
\end{align*}
Nie zakładamy $\sigma$-skończoności całej grupy.
\end{proof}

\begin{wniosek}[Zmiana postaci splotu]
\label{cor:dt-haar-convolution-change}
W przyjętej konwencji $\mu(Ea)=\Delta(a)\mu(E)$
dla $f,g\in C_c(G;\mathbb K)$ zachodzi
\[
 (f*g)(x)=\int_G f(xz^{-1})g(z)\Delta(z)^{-1}
                  \,\mathrm d\mu(z).
\]
Jeżeli grupa jest przemienna, to $f*g=g*f$.
\end{wniosek}
\begin{proof}
W definicji podstawiamy $y=xu$ (lewa translacja),
co daje
\[
 (f*g)(x)=\int_G f(xu)g(u^{-1})\,\mathrm d\mu(u).
\]
Następnie stosujemy wzór inwersji przytoczony na początku rozdziału
do funkcji $z\mapsto f(xz^{-1})g(z)$. Ma ona zwarty nośnik
zawarty w nośniku $g$, więc obie całki są skończone.
W grupie przemiennej $\Delta=1$ oraz
$xz^{-1}=z^{-1}x$, co daje przemienność splotu.
\end{proof}

\begin{przyklad}[Splot nieprzemienny]
Na skończonej nieprzemiennej grupie $G$ z miarą
liczącą wybierzmy $a,b$ takie, że $ab\ne ba$.
Wówczas $\chi_{\{a\}}*\chi_{\{b\}}=\chi_{\{ab\}}$,
ponieważ jedyny niezerowy składnik sumy odpowiada
$y=a$ i warunkowi $a^{-1}x=b$.
W odwrotnej kolejności otrzymujemy
$\chi_{\{ba\}}$, a zatem inny splot.
\end{przyklad}

\subsection{Algebra splotu i przedłużenia do przestrzeni Lebesgue'a}

\begin{twierdzenie}[Algebra splotu $L^1$]
\label{thm:dt-haar-l1-convolution}
Splot na $C_c(G;\mathbb K)$ ma dokładnie jedno
ograniczone dwuliniowe rozszerzenie do
$L^1(G,\mu;\mathbb K)\times L^1(G,\mu;\mathbb K)$
o wartościach w $L^1(G,\mu;\mathbb K)$.
Rozszerzenie jest łączne i spełnia
$\|f*g\|_1\le\|f\|_1\|g\|_1$.
Zatem $L^1(G,\mu;\mathbb K)$ ze splotem jest
algebrą Banacha, także gdy $\mu$ nie jest
$\sigma$-skończona.
\end{twierdzenie}
\begin{proof}
Miara Haara jest miarą Radona, więc z twierdzenia
\atlasref{ME-20::thm:cc-dense-lp-radon-20}{ME-20: gęstość funkcji ciągłych o zwartym nośniku} dla $p=1$ funkcje
z $C_c(G;\mathbb K)$ są gęste w $L^1(G,\mu;\mathbb K)$.
Twierdzenie to nie wymaga $\sigma$-skończoności. Ich klasy w $L^1$
jednoznacznie wyznaczają funkcje ciągłe:
jeżeli ciągła funkcja $f$ ma normę $L^1$ zero,
to dla każdego $\varepsilon>0$ otwarty zbiór
$\{|f|>\varepsilon\}$ ma miarę zero, a z wniosku
\atlasref{ME-23::cor:haar-open-positive-23}{ME-23: dodatnia miara otwartych zbiorów} musi być pusty.

Wybieramy ciągi $f_n,g_n\in C_c(G;\mathbb K)$
zbieżne w $L^1$ odpowiednio do $f,g$.
Z twierdzenia \ref{thm:dt-haar-convolution}
mamy
\[
 \|f_n*g_n-f_m*g_m\|_1
 \le\|f_n-f_m\|_1\|g_n\|_1+
    \|f_m\|_1\|g_n-g_m\|_1.
\]
Ciągi norm są ograniczone, więc prawa strona
dąży do zera przy $m,n\to\infty$. Z zupełności
$L^1$ (twierdzenie \atlasref{ME-14::thm:riesz-fischer-14}{ME-14: zupełność przestrzeni $L^1$})
istnieje granica, którą oznaczamy $f*g$.
Jeżeli $\widetilde f_n\to f$ i $\widetilde g_n\to g$
w $L^1$ dla innych ciągów z $C_c(G;\mathbb K)$, to
\[
 \|f_n*g_n-\widetilde f_n*\widetilde g_n\|_1
 \le \|f_n-\widetilde f_n\|_1\|g_n\|_1
   +\|\widetilde f_n\|_1\|g_n-\widetilde g_n\|_1
 \longrightarrow0.
\]
To dowodzi niezależności granicy od wyboru ciągów. Przechodząc do granicy
w nierówności norm, otrzymujemy
$\|f*g\|_1\le\|f\|_1\|g\|_1$.
Dwuliniowość wynika z dwuliniowości splotu
na $C_c(G)$ i liniowości granicy.
Jeśli także $h_n\in C_c(G;\mathbb K)$ i $h_n\to h$ w $L^1$,
to z otrzymanej nierówności i dwuliniowości mamy
\begin{align*}
 \|(f_n*g_n)*h_n-(f*g)*h\|_1
 &\le\|f_n*g_n-f*g\|_1\|h_n\|_1
   +\|f*g\|_1\|h_n-h\|_1,\\
 \|f_n*(g_n*h_n)-f*(g*h)\|_1
 &\le\|f_n-f\|_1\|g_n*h_n\|_1
   +\|f\|_1\|g_n*h_n-g*h\|_1.
\end{align*}
Obie prawe strony dążą do zera, ponieważ ciągi norm
$\|h_n\|_1$ oraz $\|g_n*h_n\|_1\le\|g_n\|_1\|h_n\|_1$
są ograniczone.
Równość wyrazów obu ciągów daje łączność.
Każde ograniczone dwuliniowe rozszerzenie
musi zachowywać te granice, co dowodzi
jednoznaczności. Zupełność algebry wynika
z zupełności $L^1$.

W tej ogólności definiujemy splot klas $L^1$
przez rozszerzenie ciągłe, nie przez nieuzasadnione
użycie globalnego twierdzenia Fubiniego.
\end{proof}


\begin{twierdzenie}[Nierówność Younga dla pary $L^1$ i $L^p$]
\label{thm:dt-haar-young}
Niech $1\le p<\infty$. Splot ma jednoznaczne
ograniczone dwuliniowe przedłużenie
$L^1(G,\mu)\times L^p(G,\mu)\to L^p(G,\mu)$
i spełnia
\[
 \|f*g\|_p\le \|f\|_1\|g\|_p.
\]
Dla $p=1$ jest to przedłużenie z twierdzenia
\ref{thm:dt-haar-l1-convolution}.
\end{twierdzenie}
\begin{proof}
Najpierw niech $f,g\in C_c(G;\C)$.
Dla $f=0$ teza jest oczywista. W przeciwnym razie
$\|f\|_1>0$, ponieważ niezerowa ciągła funkcja
jest oddzielona od zera na niepustym otwartym zbiorze.
Dla ustalonego $x$ nierówność Jensena względem
probabilistycznej miary $|f(y)|\,\mathrm d\mu(y)/\|f\|_1$
daje
\[
 |(f*g)(x)|^p
 \le \|f\|_1^{p-1}
 \int_G |f(y)|\,|g(y^{-1}x)|^p\,\mathrm d\mu(y).
\]
Funkcja pod całką podwójną jest ciągła i ma nośnik
w zwartym zbiorze
$\supp f\times(\supp f)(\supp g)$.
Możemy więc zastosować twierdzenie Fubiniego
dla zwartych nośników z ME-23.
Lewostronna niezmienniczość $\mu$ daje
\[
 \begin{aligned}
 \int_G |f*g|^p\,\mathrm d\mu
 &\le \|f\|_1^{p-1}\int_G|f(y)|
       \int_G |g(y^{-1}x)|^p\,\mathrm d\mu(x)
       \,\mathrm d\mu(y)\\
 &=\|f\|_1^p\|g\|_p^p.
 \end{aligned}
\]
Po wyciągnięciu pierwiastka stopnia $p$ otrzymujemy
żądaną nierówność.

Z gęstości $C_c$ w obu przestrzeniach Lebesgue'a
(por. \atlasref{ME-20::thm:cc-dense-lp-radon-20}{ME-20: gęstość})
wybieramy $f_n\to f$ w $L^1$ i $g_n\to g$ w $L^p$,
gdzie $f_n,g_n\in C_c$. Mamy
\[
 \|f_n*g_n-f_m*g_m\|_p
 \le\|f_n-f_m\|_1\|g_n\|_p+
 \|f_m\|_1\|g_n-g_m\|_p.
\]
Oba ciągi norm są ograniczone, a różnice w odpowiednich
normach dążą do zera. Z zupełności $L^p$ wynika istnienie
granicy. Dla innych ciągów przybliżających
$\widetilde f_n,\widetilde g_n$ analogiczne oszacowanie
\[
 \|f_n*g_n-\widetilde f_n*\widetilde g_n\|_p
 \le\|f_n-\widetilde f_n\|_1\|g_n\|_p+
 \|\widetilde f_n\|_1\|g_n-\widetilde g_n\|_p
\]
daje niezależność granicy. Liniowość w każdym
argumencie i nierówność norm przechodzą do granicy.
Każde ograniczone rozszerzenie musi przyjmować te
same granice, co dowodzi jednoznaczności.
Dla $p=1$ oba rozszerzenia zgadzają się na gęstej
podprzestrzeni $C_c\times C_c$, więc są równe.
Dla funkcji rzeczywistych dowód jest identyczny.
\end{proof}

\begin{uwaga}
W grupie nieunimodularnej nie wolno zamieniać ról
$f$ i $g$ w nierówności Younga bez dodatkowego
czynnika modularnego. Ponadto dla ogólnej grupy
lokalnie zwartej powyższe przedłużenia definiujemy
w normach przestrzeni Lebesgue'a, a nie przez
nieuzasadnione użycie globalnego twierdzenia Fubiniego.
\end{uwaga}
